%% =============================================================================
%%  Introduction — how to use this notebook
%% =============================================================================
\chapter{Introduction}

This notebook holds a personal study of \ebTitulo. The aim is twofold: to
understand the text as faithfully as possible and, at the same time, to
practise a \emph{method} that can be repeated in the next notebooks of the
series.

\section{The seven steps}

The path starts from what is closest to the reader (the translations already
in use) and moves towards the text in its own language, form and context,
hearing other interpreters only afterwards and then framing a synthesis. At
the end it returns to the text: interpretation is a \emph{spiral}, not a
straight line \parencite{osborne2006}.

\begin{center}
\begin{tikzpicture}[
    etapa/.style={circle, minimum size=11mm, inner sep=0pt, font=\EBftitulo\Large,
      text=white},
    rot/.style={font=\sffamily\scriptsize, text width=18mm, align=center, text=tinta},
    grupo/.style={font=\EBfversal\footnotesize, text=destaque2!80!black}]
  \foreach \i/\t/\c in {1/Translations/destaque, 2/Original text/destaque,
      3/Structure/destaque, 4/Context/grego!70!destaque2, 5/Intertext/grego!70!destaque2,
      6/Dialogue/grego!70!destaque2, 7/Synthesis/rubrica}{
    \node[etapa, fill=\c] (e\i) at ({(\i-1)*2.1}, 0) {\i};
    \node[rot, below=2pt of e\i] {\t};
  }
  \foreach \i [evaluate=\i as \j using int(\i+1)] in {1,...,6}
    \draw[-{Stealth[length=5pt]}, destaque2, thick] (e\i) -- (e\j);
  \draw[decorate, decoration={brace, amplitude=5pt}, destaque2]
    ([yshift=9mm]e1.west) -- node[grupo, above=6pt]{Observation} ([yshift=9mm]e3.east);
  \draw[decorate, decoration={brace, amplitude=5pt}, destaque2]
    ([yshift=9mm]e4.west) -- node[grupo, above=6pt]{Interpretation} ([yshift=9mm]e6.east);
  \draw[decorate, decoration={brace, amplitude=5pt}, destaque2]
    ([yshift=9mm]e7.west) -- node[grupo, above=6pt]{Appropriation} ([yshift=9mm]e7.east);
  \draw[-{Stealth[length=5pt]}, rubrica, thick, dashed, rounded corners=8pt]
    (e7.south) -- ++(0,-12mm) -| node[pos=.25, below, font=\sffamily\scriptsize, text=rubrica]
    {hermeneutical spiral: back to the text with new questions} (e1.south);
\end{tikzpicture}
\end{center}

Each step opens with a \emph{goal} box, carries numbered questions and closes
with what was learned. The steps expand a classic six-stage plan, giving the
original text, the literary structure and the use of the Old Testament a
place of their own:

\begin{ebtabela}{@{}>{\raggedright}p{.3\linewidth} >{\raggedright\arraybackslash}X@{}}
\EBcab{Classic plan} & \EBcab{In this notebook}\\\midrule
1. Translation & Step 1 — Translations (critical comparison)\\
2. Lexicon and semantics & Step 2 — Original text (textual criticism, lexicon, grammar)\\
3. Immediate context & Step 3 — Structure \emph{and} Step 4 — Context\\
4. Cross references & Step 5 — Intertext (Hebrew OT, Aramaic, LXX, Peshitta)\\
5. Commentaries & Step 6 — Dialogue with the tradition\\
6. Homiletics and doctrine & Step 7 — Synthesis and appropriation\\
\end{ebtabela}

\section{Working principles}

\paragraph{Observe before interpreting.} The first three steps only record
what is in the text: words, forms, repetitions, contrasts. Questions of
meaning are written down but left open \parencite{duvall2012, fee2002}.

\paragraph{Commentaries last.} Reading commentaries too early makes a study
repeat other people's conclusions. They come in at Step~6 as conversation
partners, to confirm, correct or extend what was observed.

\paragraph{Words mean things in context.} The interlinear is a tool for
observation, not a translation. It tends to suggest a word-for-word reading
and to hide the syntax; that is why it always goes together with an analysis
of the whole sentence.

\begin{alerta}[Fallacies to avoid]
\emph{Root fallacy}: assuming that etymology settles a word's current sense.
\emph{Illegitimate totality transfer}: loading a single occurrence with every
sense the dictionary lists. \emph{Anachronism}: reading a Greek term through
the meaning of an English word derived from it. See \textcite{carson1996} and
\textcite{silva1994}.
\end{alerta}

\section{Legend}

\noindent
\begin{tabularx}{\linewidth}{@{}l X@{}}
\gr{Ἑλληνικά} & Greek text (Gentium Book Plus, blue)\\
\hb{עִבְרִית} & Hebrew (Noto Serif Hebrew, amber)\\
\arc{אֲרָמִית} \ \sy{ܣܘܪܝܝܐ} & Aramaic and Syriac (olive green)\\
\tl{translit.} & transliteration in the SBL style\\
\vs{20}\hspace{-.2em} & verse number (rubric red, as in the manuscripts)\\
{\color{rubrica}°} & word with a recorded textual variant\\
{\color{suave}[\,]} & implied word, supplied in the gloss\\
\end{tabularx}

\begin{objetivo}States what the step sets out to achieve.\end{objetivo}
\begin{pergunta}Study question, numbered within the step.\end{pergunta}
\begin{nota}An observation, a fact or a clarification.\end{nota}
\begin{variante}A problem of textual criticism.\end{variante}
\begin{aplicacao}A bridge from the sense of the text to life today.\end{aplicacao}
